% concatenated from preprint_v3.tex
\documentclass{article}
\usepackage{arxiv}
\usepackage[utf8]{inputenc} % allow utf-8 input
\usepackage[T1]{fontenc}    % use 8-bit T1 fonts
\usepackage{hyperref}       % hyperlinks
\usepackage{url}            % simple URL typesetting
\usepackage{booktabs, rotating, amsfonts, bbm, microtype, graphicx, amsmath, amssymb, amsthm, enumitem}
\usepackage[square, sort, numbers]{natbib}
\newtheorem{assumption}{Assumption}
\newcommand{\dataset}{{\cal D}}
\newcommand{\fracpartial}[2]{\frac{\partial #1}{\partial  #2}}
\newtheorem{theorem}{Theorem}[section]
\newtheorem{lemma}[theorem]{Lemma}
\newtheorem{proposition}[theorem]{Proposition}
\newtheorem{corollary}[theorem]{Corollary}
\newtheorem{definition}[theorem]{Definition}
\newtheorem{remark}[theorem]{Remark}
\newtheorem{example}[theorem]{Example}
\numberwithin{equation}{section}
\title{A Learning Stability Profile for Finite-Dimensional Learning Dynamics}


\author{Ronald Katende \\
Department of Mathematics\\
Kabale University\\
Kikungiri Hill, Katuna Road, 317, Kabale, Uganda\\
\texttt{rkatende@kab.ac.ug} %\\
}

\date{}

\hypersetup{
pdftitle={A Learning Stability Profile for Finite-Dimensional Learning Dynamics},
pdfsubject={cs.LG, math.OC, math.DS, cs.SY},
%pdfsubject={q-bio.NC, q-bio.QM},
pdfauthor={Ronald Katende},
pdfkeywords={learning stability profile; Lyapunov stability; sensitivity analysis; nonsmooth learning systems; Clarke generalized Jacobian; stochastic gradient methods; residual networks; proximal algorithms},
}

\begin{document}
\maketitle

\begin{abstract}
We develop a finite-dimensional sensitivity framework for studying stability in learning systems whose states include representations, parameters, and update variables. The central object is the \emph{Learning Stability Profile}, a collection of directional sensitivity operators that records how perturbations in inputs, parameter initialization, and update mechanisms propagate along a specified learning trajectory.

The main result is a Lyapunov criterion for controlling this profile. Under explicit regularity, coercivity, and dissipation assumptions, an incremental Lyapunov energy yields uniform or exponentially decaying bounds on the associated linearized transition operators. The result is stated as a sufficient stability criterion, not as an unconditional converse theorem. The framework also distinguishes terminal decay, profile-wise boundedness, and subexponential growth, avoiding the identification of nonpositive growth exponents with uniform boundedness.

The profile is then specialized to several standard learning mechanisms. Spectral bounds give forward sensitivity estimates for feedforward networks. Dissipativity and step-size restrictions give stability bounds for residual architectures. Mean-square contraction assumptions yield parameter and update sensitivity bounds for stochastic gradient methods. Locally Lipschitz systems, including piecewise-linear networks, proximal maps, projected updates, and recurrent or state-space recursions, are handled through Clarke generalized Jacobians and variational Lyapunov inequalities.

The resulting framework provides a common stability language for architecture, optimization, stochasticity, and nonsmoothness. Its role is structural: it organizes known stability mechanisms within one perturbation calculus while keeping the hypotheses needed for each guarantee explicit.
\end{abstract}


% keywords can be removed
\keywords{learning stability profile \and Lyapunov stability; sensitivity analysis \and nonsmooth learning systems \and Clarke generalized Jacobian \and stochastic gradient methods \and residual networks \and proximal algorithms}

\section{Introduction}
\label{sec:introduction}

Stability enters learning theory through several distinct formalisms. In
dynamical systems, stability is commonly expressed through Lyapunov functions
and contractive evolutions \cite{HirschSmaleDevaney2012,Khalil2002Nonlinear}.
In numerical analysis, stability is often tied to step-size restrictions for
discrete schemes \cite{Strikwerda2004}. In statistical learning, algorithmic
stability measures the sensitivity of a learning rule to perturbations of the
training data \cite{Bousquet2002,Hardt2016Train}. In optimization, stability is
controlled through smoothness, monotonicity, firm nonexpansiveness, and
averagedness of update maps \cite{Nesterov2004,BauschkeCombettes2011,BottouCurtisNocedal2018}.
These notions are related, but they are usually formulated for different
objects.

This paper develops a finite-dimensional analytic framework for tracking
stability in learning systems whose states contain both representations and
parameters. The central object is the \emph{Learning Stability Profile} (LSP), a
collection of sensitivity operators describing how perturbations in inputs,
parameters, and update variables propagate along a learning trajectory. The LSP
does not replace existing stability theories. It provides a common notation in
which Lyapunov, spectral, residual, stochastic, and variational stability
estimates can be compared within one representation--parameter dynamics.

The main technical contribution is a Lyapunov criterion for controlling the
LSP. Under explicit regularity, coercivity, and dissipation assumptions, an
incremental Lyapunov energy for the linearized learning flow gives quantitative
bounds on the corresponding sensitivity operators. The result is stated as a
sufficient stability criterion. No unconditional equivalence between
nonpositive exponents and Lyapunov energies is claimed.

The framework is then specialized to standard settings. Spectral norm control
gives forward sensitivity bounds for feedforward networks. Dissipativity and a
step-size restriction give residual-network bounds analogous to stability
restrictions for explicit discretizations \cite{Strikwerda2004,HaberRuthotto2017}.
Strong convexity, smoothness, and noise bounds give mean-square estimates for
stochastic gradient methods \cite{KushnerYin2003,Benaim1999,BottouCurtisNocedal2018}.
Finally, locally Lipschitz learning maps are treated through Clarke generalized
Jacobians and variational Lyapunov functions \cite{Clarke1990Optimization,Clarke1998Nonsmooth,RockafellarWets1998}.

\section{Setup}
\label{sec:setup}

Fix a depth horizon \(L\) and a training horizon \(T\). At training time \(t\),
write the learning state as
\[
S^t
=
(X_0^t,\ldots,X_L^t,\theta^t)
\in
\mathcal{M}_L
:=
\mathcal{X}_0\times\cdots\times\mathcal{X}_L\times\Theta .
\]
Here \(X_k^t\) is the representation at depth \(k\), and \(\theta^t\) denotes the
trainable parameters. The variable \(U^t\) represents stochasticity, mini-batch
selection, update noise, or external perturbations. A training step is written as
\[
S^{t+1}
=
\Phi_t(S^t,U^t),
\qquad
t=0,\ldots,T-1 .
\]

When \(\Phi_t\) is differentiable along the trajectory, define
\[
A_t:=D_S\Phi_t(S^t,U^t),
\qquad
B_t:=D_U\Phi_t(S^t,U^t).
\]
The linearized perturbation dynamics are
\[
\delta S^{t+1}
=
A_t\delta S^t+B_t\delta U^t .
\]
For nonsmooth locally Lipschitz maps, \(A_t\) and \(B_t\) are replaced by
appropriate blocks of the Clarke generalized Jacobian
\cite{Clarke1990Optimization,Clarke1998Nonsmooth}.

The Learning Stability Profile consists of operator norms of block sensitivities
generated by the linearized system. We use
\[
\sigma^x_{k,t}
:=
\bigl\|D_{X_0^0}X_k^t\bigr\|,
\qquad
\sigma^\theta_t
:=
\bigl\|D_{\theta^0}\theta^t\bigr\|,
\qquad
\sigma^u_{s,t}
:=
\bigl\|D_{U^s}S^t\bigr\|,
\quad
0\le s<t\le T .
\]
The first quantity tracks input-to-representation sensitivity, the second
tracks sensitivity to parameter initialization, and the third tracks sensitivity
to update perturbations.

For a nonnegative sensitivity family \(\Sigma_{L,T}\), define its normalized
growth exponent by
\[
\alpha(\Sigma)
:=
\limsup_{L,T\to\infty}
\frac{1}{L+T}\log \Sigma_{L,T},
\]
whenever the expression is well defined. Terminal forward growth may be measured
by
\[
\alpha_x^{\mathrm{term}}
:=
\limsup_{L,T\to\infty}
\frac{1}{L+T}
\log
\bigl\|D_{X_0^0}X_L^T\bigr\|.
\]
Profile growth is obtained by replacing the terminal sensitivity with a maximum
over the relevant depth--time indices.

The distinction matters. A negative terminal exponent gives exponential decay
of the terminal sensitivity. A nonpositive exponent only rules out exponential
growth. It does not imply uniform boundedness. For example, polynomially growing
sensitivities have exponent zero. Uniform boundedness is therefore recorded
separately, for example
\[
\sup_{L,T}\max_{0\le k\le L,\;0\le t\le T}\sigma^x_{k,t}<\infty,
\qquad
\sup_T\max_{0\le t\le T}\sigma^\theta_t<\infty,
\qquad
\sup_T\max_{0\le s<t\le T}\sigma^u_{s,t}<\infty .
\]

We use incremental Lyapunov energies for the linearized flow. Such an energy is
a nonnegative functional \(V_t\) satisfying comparison bounds
\[
m\|\delta S\|^2
\le
V_t(\delta S)
\le
M\|\delta S\|^2
\]
for constants \(0<m\le M<\infty\), together with a discrete dissipation
inequality along the perturbation dynamics. This is the finite-dimensional
Lyapunov formulation of incremental stability \cite{Khalil2002Nonlinear,Angeli2002Incremental}.

\begin{proposition}[Basic properties of stability exponents]
\label{prop:wellposed}
Let \(\Sigma_{L,T}\ge0\). If
\[
\Sigma_{L,T}\le C
\]
uniformly in \(L,T\), then \(\alpha(\Sigma)\le0\). If
\[
\Sigma_{L,T}\le C\exp[-\lambda(L+T)]
\]
for some \(C,\lambda>0\), then \(\alpha(\Sigma)\le-\lambda\). Conversely,
\(\alpha(\Sigma)\le0\) implies only subexponential growth and does not imply
uniform boundedness. These conclusions are unchanged under norm changes whose
equivalence constants grow subexponentially in \(L+T\).
\end{proposition}

\begin{proof}
Uniform boundedness gives
\[
\frac{1}{L+T}\log \Sigma_{L,T}
\le
\frac{1}{L+T}\log C,
\]
whose limit superior is \(0\). Exponential decay gives
\[
\frac{1}{L+T}\log \Sigma_{L,T}
\le
\frac{\log C}{L+T}-\lambda,
\]
hence \(\alpha(\Sigma)\le-\lambda\). A polynomially growing sequence satisfies
\(\alpha(\Sigma)=0\) and is not uniformly bounded. If two norms differ by a
factor \(c_{L,T}\) with \(\log c_{L,T}=o(L+T)\), the normalized logarithmic
growth rates differ by \(o(1)\).
\end{proof}

\section{Analytic Lyapunov control}
\label{sec:analytic-theorem}

For a fixed trajectory \((S^t,U^t)\), consider the homogeneous linearized system
\[
\delta S^{t+1}=A_t\delta S^t .
\]
We say that this system admits a uniform incremental Lyapunov energy if there
exist \(0<m\le M<\infty\), \(0<\gamma\le M\), and functionals \(V_t\) such that
\[
m\|z\|^2
\le
V_t(z)
\le
M\|z\|^2
\]
and
\[
V_{t+1}(A_tz)-V_t(z)
\le
-\gamma\|z\|^2
\]
for all admissible perturbations \(z\).

For \(0\le s<t\), define
\[
P_{t:s}:=A_{t-1}\cdots A_s,
\qquad
P_{s:s}:=I .
\]

\begin{theorem}[Analytic Lyapunov criterion for the LSP]
\label{thm:fundamental-analytic-stability}
Assume that \(\Phi_t\) is continuously differentiable along the trajectory and
that the homogeneous linearized flow admits a uniform incremental Lyapunov
energy. Set
\[
\rho:=\left(1-\frac{\gamma}{M}\right)^{1/2}.
\]
Then \(0\le\rho<1\), and
\[
\|P_{t:s}\|
\le
\sqrt{\frac{M}{m}}\,\rho^{t-s},
\qquad
0\le s<t .
\]
Consequently, each homogeneous block sensitivity in the LSP is bounded by the
same exponential factor, up to projection constants. In particular,
\[
\sigma^\theta_t
\le
C_\theta\rho^t
\]
for a constant \(C_\theta\) depending only on the product norm and the parameter
projection.

If, in addition, \(\|B_t\|\le B\) uniformly, then
\[
\|D_{U^s}S^t\|
\le
\sqrt{\frac{M}{m}}\,B\,\rho^{t-s-1},
\qquad
0\le s<t .
\]
Thus the update-sensitivity profile is uniformly bounded:
\[
\sup_{0\le s<t}\|D_{U^s}S^t\|
\le
\sqrt{\frac{M}{m}}\,B .
\]
\end{theorem}

\begin{proof}
Let \(z_{r+1}=A_rz_r\). By the Lyapunov inequality,
\[
V_{r+1}(z_{r+1})
\le
V_r(z_r)-\gamma\|z_r\|^2 .
\]
Since \(V_r(z_r)\le M\|z_r\|^2\),
\[
\|z_r\|^2\ge \frac{1}{M}V_r(z_r).
\]
Therefore
\[
V_{r+1}(z_{r+1})
\le
\left(1-\frac{\gamma}{M}\right)V_r(z_r)
=
\rho^2 V_r(z_r).
\]
Iteration from \(s\) to \(t-1\) gives
\[
V_t(P_{t:s}z_s)
\le
\rho^{2(t-s)}V_s(z_s).
\]
Using the comparison bounds,
\[
m\|P_{t:s}z_s\|^2
\le
V_t(P_{t:s}z_s)
\le
M\rho^{2(t-s)}\|z_s\|^2 .
\]
Hence
\[
\|P_{t:s}z_s\|
\le
\sqrt{\frac{M}{m}}\,\rho^{t-s}\|z_s\|.
\]
Taking the supremum over \(z_s\ne0\) proves the transition bound.

Block sensitivity bounds follow by composing \(P_{t:s}\) with the relevant
coordinate injections and projections. For update perturbations,
\[
D_{U^s}S^t=P_{t:s+1}B_s .
\]
Thus
\[
\|D_{U^s}S^t\|
\le
\|P_{t:s+1}\|\,\|B_s\|
\le
\sqrt{\frac{M}{m}}\,B\,\rho^{t-s-1}.
\]
The uniform bound follows because \(0\le\rho<1\).
\end{proof}

The theorem is a sufficient Lyapunov criterion. It does not assert that every
subexponentially stable or nonexpansive learning system admits a strictly
dissipative energy. Converse Lyapunov statements require additional hypotheses,
such as uniform exponential stability of the linearized system
\cite{Khalil2002Nonlinear,Angeli2002Incremental}.

\section{Feedforward and residual stability laws}
\label{sec:spectral-residual-laws}

We next apply the LSP notation to standard architecture-level estimates. These
results are used as consistency checks and as sources of explicit constants.

\subsection{Feedforward networks with spectral control}
\label{subsec:feedforward-laws}

Consider
\[
X_{k+1}
=
f_k(X_k)
=
\sigma(W_kX_k+b_k),
\qquad
k=0,\ldots,L-1,
\]
where \(\sigma\) is applied componentwise and is \(1\)-Lipschitz. This includes
ReLU and leaky ReLU with slopes in \([0,1]\). Let
\[
F_\theta:=f_{L-1}\circ\cdots\circ f_0 .
\]

\begin{theorem}[Forward spectral bound]
\label{thm:ff-spectral-stability}
Assume
\[
\|W_k\|_2\le \rho_k,
\qquad
k=0,\ldots,L-1 .
\]
Then
\[
\|F_\theta(x)-F_\theta(y)\|_2
\le
\left(\prod_{k=0}^{L-1}\rho_k\right)\|x-y\|_2
\qquad
\text{for all }x,y .
\]
At differentiability points of all layers,
\[
\|JF_\theta(x)\|_2
\le
\prod_{k=0}^{L-1}\rho_k .
\]
If \(\rho_k\le\rho<1\) for all \(k\), then
\[
\|D_{X_0}X_L\|_2\le \rho^L
\]
and
\[
\alpha_x^{\mathrm{term}}
:=
\limsup_{L\to\infty}
\frac1L\log\|D_{X_0}X_L\|_2
\le
\log\rho<0 .
\]
\end{theorem}

\begin{proof}
For any \(x,y\),
\[
\|f_k(x)-f_k(y)\|_2
=
\|\sigma(W_kx+b_k)-\sigma(W_ky+b_k)\|_2 .
\]
Since \(\sigma\) is \(1\)-Lipschitz,
\[
\|f_k(x)-f_k(y)\|_2
\le
\|W_k(x-y)\|_2
\le
\rho_k\|x-y\|_2 .
\]
Composition over \(k=0,\ldots,L-1\) gives
\[
\|F_\theta(x)-F_\theta(y)\|_2
\le
\left(\prod_{k=0}^{L-1}\rho_k\right)\|x-y\|_2 .
\]

At differentiability points,
\[
Jf_k(X_k)=D_k(X_k)W_k,
\]
where \(D_k(X_k)\) is diagonal and \(\|D_k(X_k)\|_2\le1\). Hence
\[
\|Jf_k(X_k)\|_2\le \rho_k .
\]
The chain rule gives
\[
JF_\theta(x)=Jf_{L-1}(X_{L-1})\cdots Jf_0(X_0),
\]
and submultiplicativity yields
\[
\|JF_\theta(x)\|_2
\le
\prod_{k=0}^{L-1}\rho_k .
\]
If \(\rho_k\le\rho<1\), then
\[
\|D_{X_0}X_L\|_2\le\rho^L .
\]
Taking logarithms, dividing by \(L\), and passing to the limit superior gives
\[
\alpha_x^{\mathrm{term}}\le\log\rho .
\]
\end{proof}

\subsection{Residual networks and step-size stability}
\label{subsec:residual-laws}

Residual networks can be viewed as explicit depth discretizations of dynamical
systems \cite{He2016ResNet,HaberRuthotto2017,RuthottoHaber2020}. We use this
viewpoint only to derive a forward sensitivity bound.

Consider
\begin{equation}
\label{eq:resnet-layer}
X_{k+1}=X_k+h\,g_k(X_k;\theta_k),
\qquad
k=0,\ldots,L-1,
\end{equation}
where \(h>0\), and \(g_k\) is differentiable in \(x\). Write
\[
G_k(x;\theta_k):=D_xg_k(x;\theta_k).
\]
Assume that, for all admissible \((x,\theta_k)\),
\begin{equation}
\label{eq:resnet-dissipative-assumption}
\frac{G_k(x;\theta_k)+G_k(x;\theta_k)^\top}{2}
\preceq
-mI,
\qquad
\|G_k(x;\theta_k)\|_2\le M_g ,
\end{equation}
with \(m>0\) and \(M_g<\infty\). These are standard dissipativity and bounded
Jacobian hypotheses for explicit stability estimates \cite{Strikwerda2004,Khalil2002Nonlinear}.

\begin{lemma}[One-step residual sensitivity]
\label{lem:resnet-one-step}
Under \eqref{eq:resnet-dissipative-assumption},
\[
\|I+hG_k(x;\theta_k)\|_2^2
\le
1-2hm+h^2M_g^2 .
\]
\end{lemma}

\begin{proof}
Let \(\|v\|_2=1\). Then
\[
\|(I+hG_k)v\|_2^2
=
1+2h\,v^\top G_kv+h^2\|G_kv\|_2^2 .
\]
Since
\[
v^\top G_kv
=
v^\top\frac{G_k+G_k^\top}{2}v
\le
-m
\]
and \(\|G_kv\|_2^2\le M_g^2\), we obtain
\[
\|(I+hG_k)v\|_2^2
\le
1-2hm+h^2M_g^2 .
\]
Taking the supremum over unit \(v\) proves the claim.
\end{proof}

Define
\[
q_h:=1-2hm+h^2M_g^2,
\qquad
c_h:=q_h^{1/2}.
\]

\begin{theorem}[Forward residual sensitivity bound]
\label{thm:resnet-analytic-stability}
Assume \eqref{eq:resnet-dissipative-assumption} and
\[
0<h<\frac{2m}{M_g^2}.
\]
Then \(0\le q_h<1\), and
\[
\|D_{X_0}X_L\|_2\le c_h^L .
\]
Consequently,
\[
\alpha_x^{\mathrm{term}}
\le
\log c_h
=
\frac12\log(1-2hm+h^2M_g^2)
<0,
\]
with the convention that the exponent is \(-\infty\) if \(c_h=0\).
\end{theorem}

\begin{proof}
Differentiating \eqref{eq:resnet-layer} gives
\[
D_{X_k}X_{k+1}=I+hG_k(X_k;\theta_k).
\]
By Lemma~\ref{lem:resnet-one-step},
\[
\|D_{X_k}X_{k+1}\|_2\le c_h .
\]
The step-size condition gives \(q_h<1\). Since \(q_h\) bounds a squared norm from
above, the assumptions are consistent only when \(q_h\ge0\). Hence
\(0\le c_h<1\).

The chain rule gives
\[
D_{X_0}X_L
=
D_{X_{L-1}}X_L\cdots D_{X_0}X_1 .
\]
Therefore
\[
\|D_{X_0}X_L\|_2
\le
\prod_{k=0}^{L-1}\|D_{X_k}X_{k+1}\|_2
\le
c_h^L .
\]
Taking logarithms, dividing by \(L\), and passing to the limit superior gives
\[
\alpha_x^{\mathrm{term}}\le\log c_h .
\]
\end{proof}

\section{Temporal stability of stochastic gradient methods}
\label{sec:sgd-temporal}

Consider
\begin{equation}
\label{eq:sgd-update}
\theta^{t+1}
=
\theta^t-\eta_tG(\theta^t,u^t),
\end{equation}
where \(u^t\) represents update randomness. The following assumptions are
standard in mean-square analyses of stochastic approximation
\cite{KushnerYin2003,Benaim1999,BottouCurtisNocedal2018}.

\begin{assumption}[Stochastic gradient regime]
\label{ass:sgd-regime}
The objective \(\mathcal{L}\) is \(\mu\)-strongly convex and \(L\)-smooth, with
minimizer \(\theta^\star\). The stochastic gradient satisfies
\[
\mathbb{E}[G(\theta,u)\mid\theta]
=
\nabla\mathcal{L}(\theta),
\]
and
\[
\mathbb{E}
\left[
\|G(\theta,u)-\nabla\mathcal{L}(\theta)\|^2
\mid\theta
\right]
\le
\sigma_0^2+\sigma_1^2\|\theta-\theta^\star\|^2 .
\]
Moreover, for some \(L_G,L_u^G\ge0\),
\begin{equation}
\label{eq:sgd-Lpar-Lupd}
\|G(\theta_1,u)-G(\theta_2,u)\|
\le
L_G\|\theta_1-\theta_2\|,
\qquad
\|G(\theta,u_1)-G(\theta,u_2)\|
\le
L_u^G\|u_1-u_2\|.
\end{equation}
\end{assumption}

Let
\[
\Delta^t:=\theta^t-\theta^\star .
\]

\begin{lemma}[Mean-square Lyapunov recursion]
\label{lem:sgd-lyapunov}
Under Assumption~\ref{ass:sgd-regime},
\[
\mathbb{E}
\left[
\|\Delta^{t+1}\|^2
\mid
\theta^t
\right]
\le
\left(1-2\mu\eta_t+\eta_t^2(L^2+\sigma_1^2)\right)\|\Delta^t\|^2
+
\eta_t^2\sigma_0^2 .
\]
\end{lemma}

\begin{proof}
From \eqref{eq:sgd-update},
\[
\Delta^{t+1}
=
\Delta^t-\eta_tG(\theta^t,u^t).
\]
Thus
\[
\|\Delta^{t+1}\|^2
=
\|\Delta^t\|^2
-
2\eta_t\langle \Delta^t,G(\theta^t,u^t)\rangle
+
\eta_t^2\|G(\theta^t,u^t)\|^2 .
\]
Taking conditional expectation gives
\[
\mathbb{E}
[
\|\Delta^{t+1}\|^2
\mid
\theta^t
]
=
\|\Delta^t\|^2
-
2\eta_t
\langle
\Delta^t,\nabla\mathcal{L}(\theta^t)
\rangle
+
\eta_t^2
\mathbb{E}
[
\|G(\theta^t,u^t)\|^2
\mid
\theta^t
].
\]
Strong convexity and \(\nabla\mathcal{L}(\theta^\star)=0\) imply
\[
\langle \Delta^t,\nabla\mathcal{L}(\theta^t)\rangle
\ge
\mu\|\Delta^t\|^2 .
\]
Let
\[
\xi^t:=G(\theta^t,u^t)-\nabla\mathcal{L}(\theta^t).
\]
By unbiasedness, \(\mathbb{E}[\xi^t\mid\theta^t]=0\), and hence
\[
\mathbb{E}
[
\|G(\theta^t,u^t)\|^2
\mid
\theta^t
]
=
\|\nabla\mathcal{L}(\theta^t)\|^2
+
\mathbb{E}
[
\|\xi^t\|^2
\mid
\theta^t
].
\]
Smoothness gives
\[
\|\nabla\mathcal{L}(\theta^t)\|
\le
L\|\Delta^t\|.
\]
The variance bound then gives
\[
\mathbb{E}
[
\|G(\theta^t,u^t)\|^2
\mid
\theta^t
]
\le
(L^2+\sigma_1^2)\|\Delta^t\|^2+\sigma_0^2 .
\]
Substitution proves the claim.
\end{proof}

\subsection{Constant and decreasing step sizes}
\label{subsec:sgd-temporal-stability}

The previous lemma controls optimization error. LSP sensitivities require a
separate contraction condition for coupled perturbations.

For \(\eta>0\), define
\[
T_\eta(\theta,u):=\theta-\eta G(\theta,u).
\]

\begin{assumption}[Mean-square sensitivity contraction]
\label{ass:sgd-sensitivity-contraction}
For the step sizes considered below, there exist \(\kappa_t\in[0,1]\) such that
coupled runs with the same update randomness satisfy
\[
\mathbb{E}
\left[
\|T_{\eta_t}(\theta_1,u^t)-T_{\eta_t}(\theta_2,u^t)\|^2
\,\middle|\,
\theta_1,\theta_2
\right]
\le
\kappa_t^2\|\theta_1-\theta_2\|^2 .
\]
For constant step size, write \(\kappa_t=\kappa_\eta<1\).
\end{assumption}

This assumption is the sensitivity counterpart of contractivity of the update
map. For instance, if the sample-gradient maps are uniformly strongly monotone
and Lipschitz, the usual gradient-step calculation gives such a contraction
under the corresponding step-size restriction \cite{Nesterov2004,BottouCurtisNocedal2018}.

\begin{theorem}[Constant-step mean-square stability]
\label{thm:sgd-constant-stepsize}
Let Assumption~\ref{ass:sgd-regime} hold and let \(\eta_t\equiv\eta\). Define
\[
q_\eta:=1-2\mu\eta+\eta^2(L^2+\sigma_1^2).
\]
If
\[
0<\eta<\frac{2\mu}{L^2+\sigma_1^2},
\]
then \(0\le q_\eta<1\) and
\[
\mathbb{E}\|\Delta^t\|^2
\le
q_\eta^t\|\Delta^0\|^2
+
\eta^2\sigma_0^2\frac{1-q_\eta^t}{1-q_\eta}.
\]
If Assumption~\ref{ass:sgd-sensitivity-contraction} also holds with
\(\kappa_\eta<1\), then
\[
\left(
\mathbb{E}\|D_{\theta^0}\theta^t\|^2
\right)^{1/2}
\le
\kappa_\eta^t ,
\]
and
\[
\alpha_{\theta,\mathrm{ms}}^{\mathrm{term}}
:=
\limsup_{t\to\infty}
\frac1t
\log
\left(
\mathbb{E}\|D_{\theta^0}\theta^t\|^2
\right)^{1/2}
\le
\log\kappa_\eta<0 .
\]
For \(0\le s<t\),
\[
\left(
\mathbb{E}\|D_{u^s}\theta^t\|^2
\right)^{1/2}
\le
\eta L_u^G\kappa_\eta^{t-s-1}.
\]
Consequently,
\[
\sup_{0\le s<t}
\left(
\mathbb{E}\|D_{u^s}\theta^t\|^2
\right)^{1/2}
\le
\eta L_u^G .
\]
\end{theorem}

\begin{proof}
Taking expectation in Lemma~\ref{lem:sgd-lyapunov} gives
\[
\mathbb{E}\|\Delta^{t+1}\|^2
\le
q_\eta\mathbb{E}\|\Delta^t\|^2+\eta^2\sigma_0^2 .
\]
Iteration yields
\[
\mathbb{E}\|\Delta^t\|^2
\le
q_\eta^t\|\Delta^0\|^2
+
\eta^2\sigma_0^2\sum_{j=0}^{t-1}q_\eta^j
=
q_\eta^t\|\Delta^0\|^2
+
\eta^2\sigma_0^2\frac{1-q_\eta^t}{1-q_\eta}.
\]

Let two coupled trajectories be driven by the same \(u^t\), and set
\[
\delta^t:=\theta_1^t-\theta_2^t .
\]
Assumption~\ref{ass:sgd-sensitivity-contraction} gives
\[
\mathbb{E}[\|\delta^{t+1}\|^2\mid \delta^t]
\le
\kappa_\eta^2\|\delta^t\|^2 .
\]
Thus
\[
\mathbb{E}\|\delta^t\|^2
\le
\kappa_\eta^{2t}\|\delta^0\|^2 .
\]
Taking \(\delta^0\) as a unit perturbation in \(\theta^0\) gives
\[
\left(
\mathbb{E}\|D_{\theta^0}\theta^t\|^2
\right)^{1/2}
\le
\kappa_\eta^t .
\]
The exponent bound follows by taking normalized logarithms.

For update perturbations,
\[
\|D_{u^s}\theta^{s+1}\|\le \eta L_u^G .
\]
Propagation from \(s+1\) to \(t\) is bounded by
\(\kappa_\eta^{t-s-1}\) in mean square. Hence
\[
\left(
\mathbb{E}\|D_{u^s}\theta^t\|^2
\right)^{1/2}
\le
\eta L_u^G\kappa_\eta^{t-s-1}.
\]
Taking the supremum over \(s<t\) proves the final bound.
\end{proof}

\begin{theorem}[Decreasing-step sensitivity bounds]
\label{thm:sgd-decreasing-stepsize}
Let Assumption~\ref{ass:sgd-regime} hold and suppose
\[
\eta_t>0,
\qquad
\sum_{t=0}^{\infty}\eta_t=\infty,
\qquad
\sum_{t=0}^{\infty}\eta_t^2<\infty .
\]
Under the standard boundedness hypotheses of stochastic approximation,
\[
\theta^t\to\theta^\star
\quad
\text{almost surely}
\]
\cite{KushnerYin2003,Benaim1999}. If Assumption~\ref{ass:sgd-sensitivity-contraction}
holds with
\[
\prod_{r=0}^{t-1}\kappa_r\to0,
\qquad
\sup_t\kappa_t\le1,
\]
then
\[
\left(
\mathbb{E}\|D_{\theta^0}\theta^t\|^2
\right)^{1/2}
\le
\prod_{r=0}^{t-1}\kappa_r ,
\qquad
\alpha_{\theta,\mathrm{ms}}^{\mathrm{term}}\le0 .
\]
For \(0\le s<t\),
\[
\left(
\mathbb{E}\|D_{u^s}\theta^t\|^2
\right)^{1/2}
\le
\eta_sL_u^G
\prod_{r=s+1}^{t-1}\kappa_r .
\]
\end{theorem}

\begin{proof}
The almost-sure convergence statement is the classical stochastic approximation
conclusion under the stated step-size conditions and boundedness assumptions
\cite{KushnerYin2003,Benaim1999}.

For sensitivities, recursive application of
Assumption~\ref{ass:sgd-sensitivity-contraction} gives
\[
\mathbb{E}\|\delta^t\|^2
\le
\left(\prod_{r=0}^{t-1}\kappa_r^2\right)\|\delta^0\|^2 .
\]
Taking \(\delta^0\) as a unit parameter perturbation gives
\[
\left(
\mathbb{E}\|D_{\theta^0}\theta^t\|^2
\right)^{1/2}
\le
\prod_{r=0}^{t-1}\kappa_r .
\]
Since the product is at most one, the normalized logarithmic growth is
nonpositive.

For update perturbations,
\[
\|D_{u^s}\theta^{s+1}\|\le \eta_sL_u^G .
\]
Propagation from \(s+1\) to \(t\) gives
\[
\left(
\mathbb{E}\|D_{u^s}\theta^t\|^2
\right)^{1/2}
\le
\eta_sL_u^G
\prod_{r=s+1}^{t-1}\kappa_r .
\]
\end{proof}

Constant step size gives a noise floor in the optimization error. Sensitivity
to initialization can still decay exponentially when the coupled update map is
contractive. Decreasing step sizes generally yield nonpositive, not necessarily
negative, normalized sensitivity exponents.

\section{Variational extension to nonsmooth learning systems}
\label{sec:variational-extension}

We now extend the LSP to locally Lipschitz learning maps. This covers ReLU-type
activations, nonsmooth losses, clipping, and quantisation through Clarke
generalized Jacobians \cite{Clarke1990Optimization,Clarke1998Nonsmooth}. It also
covers proximal maps and projections, which are nonexpansive in the convex
setting \cite{BauschkeCombettes2011,RockafellarWets1998}.

\subsection{Locally Lipschitz learning maps}
\label{subsec:locally-lipschitz-learning-maps}

Write one learning step as
\[
S^{t+1}=\Psi_t(S^t,U^t),
\qquad
S^t\in\mathcal{M}_L .
\]

\begin{assumption}[Locally Lipschitz learning flow]
\label{ass:LL}
For each \(t\), the map
\[
\Psi_t:\mathcal{M}_L\times\mathcal{U}\to\mathcal{M}_L
\]
is locally Lipschitz near \((S^t,U^t)\). For every
\[
H_t\in\partial_C\Psi_t(S^t,U^t),
\]
write
\[
H_t(\delta S,\delta U)=V_t\delta S+W_t\delta U .
\]
There exist \(L_S^C,L_U^C<\infty\) such that
\[
\|V_t\|\le L_S^C,
\qquad
\|W_t\|\le L_U^C
\]
for all admissible \(H_t\) along the trajectory.
\end{assumption}

In finite dimensions, locally Lipschitz maps have nonempty compact convex
Clarke generalized Jacobians \cite{Clarke1990Optimization}. The generalized
perturbation dynamics are represented by
\[
\delta S^{t+1}\in V_t\delta S^t+W_t\delta U^t,
\qquad
H_t=(V_t,W_t)\in\partial_C\Psi_t(S^t,U^t).
\]

For \(0\le s<t\), define
\[
\mathcal{P}_{t:s}^C
:=
\left\{
V_{t-1}\cdots V_s:
H_r=(V_r,W_r)\in\partial_C\Psi_r(S^r,U^r),
\ r=s,\ldots,t-1
\right\},
\qquad
\mathcal{P}_{s:s}^C:=\{I\}.
\]
Let \(\Pi_x,\Pi_\theta\) be the representation and parameter projections, and
let \(\iota_x,\iota_\theta\) be the corresponding injections. Define
\[
\sigma_{x,t}^C
:=
\sup_{P\in\mathcal{P}_{t:0}^C}
\|\Pi_xP\iota_x\|,
\qquad
\sigma_{\theta,t}^C
:=
\sup_{P\in\mathcal{P}_{t:0}^C}
\|\Pi_\theta P\iota_\theta\|.
\]
For \(0\le s<t\), define
\[
\sigma_{u,s,t}^C
:=
\sup
\left\{
\|PW_s\|:
P\in\mathcal{P}_{t:s+1}^C,\
H_s=(V_s,W_s)\in\partial_C\Psi_s(S^s,U^s)
\right\}.
\]
If \(\Psi_t\) is differentiable, these quantities reduce to the smooth LSP.

The terminal generalized exponents are
\[
\tilde\alpha_x^{\mathrm{term}}
:=
\limsup_{t\to\infty}
\frac1t\log\sigma_{x,t}^C,
\qquad
\tilde\alpha_\theta^{\mathrm{term}}
:=
\limsup_{t\to\infty}
\frac1t\log\sigma_{\theta,t}^C .
\]
The update-profile exponent is
\[
\tilde\alpha_u^{\mathrm{prof}}
:=
\limsup_{T\to\infty}
\frac1T
\log
\sup_{0\le s<t\le T}\sigma_{u,s,t}^C .
\]
As before, negative terminal exponents imply exponential decay, whereas
nonpositive profile exponents only rule out exponential growth.

\subsection{Variational Lyapunov control}
\label{subsec:variational-lyapunov-control}

\begin{assumption}[Incremental variational Lyapunov energy]
\label{ass:var-energy-main}
There exist \(0<m\le M<\infty\), \(0<\gamma\le M\), and proper lower
semicontinuous functionals
\[
\mathcal{V}_t:\mathcal{M}_L\to\mathbb{R}_+
\]
such that, for every admissible \(z\in\mathcal{M}_L\),
\[
m\|z\|^2
\le
\mathcal{V}_t(z)
\le
M\|z\|^2 ,
\]
and, for every
\[
H_t=(V_t,W_t)\in\partial_C\Psi_t(S^t,U^t),
\]
\[
\mathcal{V}_{t+1}(V_tz)-\mathcal{V}_t(z)
\le
-\gamma\|z\|^2 .
\]
\end{assumption}

\begin{theorem}[Variational Lyapunov bound for the generalized LSP]
\label{thm:fundamental-variational-main}
Let Assumptions~\ref{ass:LL} and~\ref{ass:var-energy-main} hold. Set
\[
\rho:=\left(1-\frac{\gamma}{M}\right)^{1/2}.
\]
Then \(0\le\rho<1\), and every \(P\in\mathcal{P}_{t:s}^C\) satisfies
\[
\|P\|
\le
\sqrt{\frac{M}{m}}\,\rho^{t-s}.
\]
Consequently, there exist constants \(C_x,C_\theta>0\), depending only on
\((m,M)\) and the relevant projections and injections, such that
\[
\sigma_{x,t}^C\le C_x\rho^t,
\qquad
\sigma_{\theta,t}^C\le C_\theta\rho^t .
\]
Thus
\[
\tilde\alpha_x^{\mathrm{term}}\le\log\rho<0,
\qquad
\tilde\alpha_\theta^{\mathrm{term}}\le\log\rho<0 .
\]
Moreover,
\[
\sigma_{u,s,t}^C
\le
\sqrt{\frac{M}{m}}\,L_U^C\,\rho^{t-s-1},
\qquad
0\le s<t,
\]
and
\[
\sup_{0\le s<t}\sigma_{u,s,t}^C
\le
\sqrt{\frac{M}{m}}\,L_U^C .
\]
\end{theorem}

\begin{proof}
Let
\[
P=V_{t-1}\cdots V_s\in\mathcal{P}_{t:s}^C .
\]
For \(z_s\in\mathcal{M}_L\), set
\[
z_{r+1}:=V_rz_r,
\qquad
r=s,\ldots,t-1 .
\]
By Assumption~\ref{ass:var-energy-main},
\[
\mathcal{V}_{r+1}(z_{r+1})
\le
\mathcal{V}_r(z_r)-\gamma\|z_r\|^2 .
\]
Since \(\mathcal{V}_r(z_r)\le M\|z_r\|^2\),
\[
\|z_r\|^2\ge \frac{1}{M}\mathcal{V}_r(z_r).
\]
Therefore
\[
\mathcal{V}_{r+1}(z_{r+1})
\le
\left(1-\frac{\gamma}{M}\right)\mathcal{V}_r(z_r)
=
\rho^2\mathcal{V}_r(z_r).
\]
Iteration gives
\[
\mathcal{V}_t(Pz_s)
\le
\rho^{2(t-s)}\mathcal{V}_s(z_s).
\]
Using the comparison bounds,
\[
m\|Pz_s\|^2
\le
\mathcal{V}_t(Pz_s)
\le
M\rho^{2(t-s)}\|z_s\|^2 .
\]
Hence
\[
\|Pz_s\|
\le
\sqrt{\frac{M}{m}}\rho^{t-s}\|z_s\|.
\]
Taking the supremum over \(z_s\ne0\) gives
\[
\|P\|
\le
\sqrt{\frac{M}{m}}\rho^{t-s}.
\]

For the representation component,
\[
\|\Pi_xP\iota_x\|
\le
\|\Pi_x\|\,\|P\|\,\|\iota_x\|
\le
\|\Pi_x\|\,\|\iota_x\|\sqrt{\frac{M}{m}}\rho^t .
\]
Thus
\[
\sigma_{x,t}^C\le C_x\rho^t,
\qquad
C_x:=\|\Pi_x\|\,\|\iota_x\|\sqrt{\frac{M}{m}} .
\]
The parameter bound follows with
\[
C_\theta:=\|\Pi_\theta\|\,\|\iota_\theta\|\sqrt{\frac{M}{m}} .
\]
The terminal exponent bounds follow by taking logarithms, dividing by \(t\),
and passing to the limit superior.

For update sensitivity, a perturbation at time \(s\) enters through some
\(W_s\) and is propagated by an element of \(\mathcal{P}_{t:s+1}^C\). Hence
\[
\|PW_s\|
\le
\|P\|\,\|W_s\|
\le
\sqrt{\frac{M}{m}}\,L_U^C\,\rho^{t-s-1}.
\]
Taking the supremum over admissible \(P\) and \(W_s\) gives the stated bound for
\(\sigma_{u,s,t}^C\). Since \(0\le\rho<1\), the profile is uniformly bounded.
\end{proof}

Theorem~\ref{thm:fundamental-variational-main} is a one-way Lyapunov criterion, i.e,. uniform variational dissipation controls the generalized LSP. No converse is claimed. A converse would require additional hypotheses, such as uniform exponential stability of the associated difference inclusion and the assumptions of a specific nonsmooth converse Lyapunov theorem.

\section{Scope of the stability profile}
\label{sec:scope-stability-profile}

The Learning Stability Profile is a finite-dimensional sensitivity framework. It does not assert a universal stability theorem for all learning systems. Its role is to organize several stability mechanisms through the same perturbation calculus: homogeneous sensitivity to initialization, forced sensitivity to update perturbations, terminal growth exponents, and profile-wise boundedness.

The framework covers three levels of stability.

First, it covers finite-horizon sensitivity bounds. These are the primary objects in the theory, since practical learning systems are trained and evaluated over finite depth and finite time. Bounds of the form
\[
\|D_{z_0}z_t\|\le C\rho^t
\quad\text{or}\quad
\sup_{0\le s<t\le T}\|D_{u^s}z_t\|\le C
\] give direct control of perturbation amplification over the learning trajectory.

Second, it covers asymptotic growth rates through normalized logarithmic exponents. Negative exponents imply exponential decay of the corresponding terminal sensitivities. Nonpositive exponents rule out exponential growth but do not imply uniform boundedness. Uniform boundedness is therefore treated as a separate estimate.

Third, it covers Lyapunov-controlled regimes. In smooth systems, a coercive incremental Lyapunov energy gives explicit bounds on the linearized transition operators. In locally Lipschitz systems, the same argument applies to admissible Clarke generalized transition products. These are sufficient stability criteria, not unconditional converse statements. Converse Lyapunov theorems require additional assumptions, such as uniform exponential stability of the linearized flow or of the associated difference inclusion \cite{Khalil2002Nonlinear,Angeli2002Incremental,Clarke1998Nonsmooth}.

The framework is intentionally local to a specified learning trajectory unless global assumptions are imposed. Global spectral bounds, global dissipativity, or global Lipschitz constants yield global stability estimates. Local versions of the same assumptions yield local or trajectory-dependent estimates. This distinction is essential for large learning systems, where global constants may be too conservative or may fail to exist.

The framework also separates deterministic, stochastic, and variational notions of stability. Deterministic systems are controlled through operator norms of the linearized flow. Stochastic systems require the mode of control to be specified, for example almost sure, in expectation, in mean square, or with high probability. Nonsmooth systems require generalized derivatives or variational transition sets. These choices change the meaning of the stability bound, but not the underlying profile structure.

Several regimes are outside the present claims. The paper does not give sharp bounds for arbitrary adaptive optimizers, nonstationary data streams, heavy-tailed noise without moment control, unbounded parameter trajectories, or infinite-dimensional limits. Such cases may still be studied using the profile, but they require additional assumptions and separate estimates. The contribution of the present framework is therefore structural, i.e., it gives a common sensitivity language and Lyapunov route for stability analysis, while keeping the hypotheses needed for each stability guarantee explicit.

\section{Worked examples}
\label{sec:worked-examples}

We give three examples showing how the Learning Stability Profile is computed in
common learning mechanisms. The first concerns nonsmooth architectures, the
second variational optimization, and the third recurrent or state-space
propagation. The examples are not presented as new stability theorems for these
individual models. Their role is to show that the same sensitivity notation
captures different stability mechanisms without changing the underlying
calculus.

\subsection{Piecewise-linear networks and active-set stability}
\label{subsec:piecewise-linear-active-set}

Consider a feedforward network
\[
X_{k+1}
=
f_k(X_k)
=
\sigma(W_kX_k+b_k),
\qquad
k=0,\ldots,L-1,
\]
where \(\sigma\) is applied componentwise and is piecewise affine with slopes in
\([0,1]\). This includes ReLU and leaky ReLU with slopes in \([0,1]\). Such maps
are locally Lipschitz, and their Clarke generalized Jacobians are well defined
\cite{Clarke1990Optimization,Clarke1998Nonsmooth}.

For \(a_k(x):=W_kx+b_k\), define the active-set diagonal family
\[
\mathcal{D}_k(x)
:=
\left\{
D=\operatorname{diag}(d_1,\ldots,d_{n_k}):
d_i\in \partial_C\sigma((a_k(x))_i)
\right\}.
\]
Then
\[
\partial_C f_k(x)
\subseteq
\{DW_k:\ D\in\mathcal{D}_k(x)\}.
\]
Since \(\partial_C\sigma(r)\subseteq[0,1]\), each \(D\in\mathcal{D}_k(x)\)
satisfies \(\|D\|_2\le1\).

\begin{proposition}[Active-set forward bound]
\label{prop:active-set-forward-bound}
Let
\[
F_\theta:=f_{L-1}\circ\cdots\circ f_0 .
\]
For every admissible generalized transition product
\[
P_L
=
D_{L-1}W_{L-1}\cdots D_0W_0,
\qquad
D_k\in\mathcal{D}_k(X_k),
\]
one has
\[
\|P_L\|_2
\le
\prod_{k=0}^{L-1}\|W_k\|_2 .
\]
Consequently, if \(\|W_k\|_2\le\rho<1\) for all \(k\), then
\[
\sup_{P_L}\|P_L\|_2
\le
\rho^L
\]
and
\[
\tilde\alpha_x^{\mathrm{term}}
\le
\log\rho<0 .
\]
\end{proposition}

\begin{proof}
For each admissible active-set matrix \(D_k\),
\[
\|D_kW_k\|_2
\le
\|D_k\|_2\|W_k\|_2
\le
\|W_k\|_2 .
\]
Hence, by submultiplicativity,
\[
\|P_L\|_2
=
\|D_{L-1}W_{L-1}\cdots D_0W_0\|_2
\le
\prod_{k=0}^{L-1}\|D_kW_k\|_2
\le
\prod_{k=0}^{L-1}\|W_k\|_2 .
\]
If \(\|W_k\|_2\le\rho<1\), then
\[
\sup_{P_L}\|P_L\|_2
\le
\rho^L .
\]
Taking logarithms, dividing by \(L\), and passing to the limit superior gives
\[
\tilde\alpha_x^{\mathrm{term}}
\le
\limsup_{L\to\infty}\frac1L\log(\rho^L)
=
\log\rho<0 .
\]
\end{proof}

The same estimate yields an explicit variational dissipation along depth. If
\[
\mathcal{V}_k(z):=\|z\|_2^2
\]
and \(\|W_k\|_2\le\rho<1\), then for every admissible \(D_k\),
\[
\mathcal{V}_{k+1}(D_kW_kz)-\mathcal{V}_k(z)
=
\|D_kW_kz\|_2^2-\|z\|_2^2
\le
-(1-\rho^2)\|z\|_2^2 .
\]
Thus the variational Lyapunov condition holds with \(m=M=1\) and
\(\gamma=1-\rho^2\). This gives the same exponent bound through
Theorem~\ref{thm:fundamental-variational-main}.

\subsection{Proximal and projected gradient flows}
\label{subsec:proximal-projected-example}

Let
\[
\mathcal{F}(\theta):=\mathcal{L}(\theta)+R(\theta),
\]
where \(R\) is proper, closed, and convex, and \(\mathcal{L}\in C^2\) satisfies
\[
\mu I\preceq \nabla^2\mathcal{L}(\theta)\preceq LI
\qquad
\text{for all }\theta ,
\]
with \(0<\mu\le L<\infty\). Consider the proximal gradient map
\[
T_\eta(\theta)
:=
\operatorname{prox}_{\eta R}
\bigl(\theta-\eta\nabla\mathcal{L}(\theta)\bigr).
\]
The proximal map of a closed convex function is firmly nonexpansive
\cite{BauschkeCombettes2011}. If \(R\) is the indicator function of a closed
convex set, \(T_\eta\) becomes projected gradient descent
\cite{RockafellarWets1998}.

\begin{proposition}[Proximal sensitivity bound]
\label{prop:proximal-sensitivity-bound}
Let
\[
0<\eta<\frac{2}{L+\mu}
\]
and define
\[
q_\eta
:=
\max\{|1-\eta\mu|,\ |1-\eta L|\}.
\]
Then \(0\le q_\eta<1\) and
\[
\|T_\eta(\theta_1)-T_\eta(\theta_2)\|
\le
q_\eta\|\theta_1-\theta_2\|
\qquad
\text{for all }\theta_1,\theta_2 .
\]
Consequently, the iterates \(\theta^{t+1}=T_\eta(\theta^t)\) satisfy
\[
\|D_{\theta^0}\theta^t\|
\le
q_\eta^t
\]
at differentiability points, and the terminal parametric exponent obeys
\[
\alpha_\theta^{\mathrm{term}}
\le
\log q_\eta<0 .
\]
\end{proposition}

\begin{proof}
For any \(\theta_1,\theta_2\), write
\[
d:=\theta_1-\theta_2 .
\]
By the mean value formula,
\[
\nabla\mathcal{L}(\theta_1)-\nabla\mathcal{L}(\theta_2)
=
H_{\theta_1,\theta_2}d,
\]
where
\[
H_{\theta_1,\theta_2}
:=
\int_0^1
\nabla^2\mathcal{L}\bigl(\theta_2+s(\theta_1-\theta_2)\bigr)\,ds .
\]
The Hessian bounds imply
\[
\mu I\preceq H_{\theta_1,\theta_2}\preceq LI .
\]
Hence
\[
\left\|
d-\eta\bigl(\nabla\mathcal{L}(\theta_1)-\nabla\mathcal{L}(\theta_2)\bigr)
\right\|
=
\|(I-\eta H_{\theta_1,\theta_2})d\|
\le
q_\eta\|d\|.
\]
Since \(\operatorname{prox}_{\eta R}\) is nonexpansive,
\[
\|T_\eta(\theta_1)-T_\eta(\theta_2)\|
\le
q_\eta\|\theta_1-\theta_2\|.
\]
The condition \(0<\eta<2/(L+\mu)\) gives \(q_\eta<1\). Iterating the contraction
gives
\[
\|D_{\theta^0}\theta^t\|
\le
q_\eta^t ,
\]
and the exponent bound follows by taking normalized logarithms.
\end{proof}

The same map has a variational descent estimate. Let
\[
\theta^+:=T_\eta(\theta).
\]
For \(0<\eta<1/L\),
\[
\mathcal{F}(\theta^+)-\mathcal{F}(\theta)
\le
-\left(\frac{1}{2\eta}-\frac{L}{2}\right)\|\theta^+-\theta\|^2 .
\]

Indeed, the proximal optimality inequality gives
\[
R(\theta^+)
+
\frac{1}{2\eta}
\|\theta^+-\theta+\eta\nabla\mathcal{L}(\theta)\|^2
\le
R(\theta)
+
\frac{\eta}{2}\|\nabla\mathcal{L}(\theta)\|^2 .
\]
Expanding and cancelling the common term
\(\frac{\eta}{2}\|\nabla\mathcal{L}(\theta)\|^2\) gives
\[
R(\theta^+)
\le
R(\theta)
-
\langle\nabla\mathcal{L}(\theta),\theta^+-\theta\rangle
-
\frac{1}{2\eta}\|\theta^+-\theta\|^2 .
\]
Smoothness gives
\[
\mathcal{L}(\theta^+)
\le
\mathcal{L}(\theta)
+
\langle\nabla\mathcal{L}(\theta),\theta^+-\theta\rangle
+
\frac{L}{2}\|\theta^+-\theta\|^2 .
\]
Adding the two inequalities yields
\[
\mathcal{F}(\theta^+)-\mathcal{F}(\theta)
\le
-\left(\frac{1}{2\eta}-\frac{L}{2}\right)\|\theta^+-\theta\|^2 .
\]
Thus proximal and projected flows fit both the sensitivity and variational sides of the LSP.

\subsection{Recurrent and state-space propagation}
\label{subsec:recurrent-state-space-example}

Consider a recurrent or state-space update
\[
H_{t+1}
=
\varphi(A_tH_t+B_tR_t+b_t),
\qquad
t=0,\ldots,T-1,
\]
where \(H_t\) is the hidden state, \(R_t\) is an external input or update signal, and \(\varphi\) is componentwise \(1\)-Lipschitz. Such recursions are standard finite-dimensional state-space systems; stability is governed by sensitivity of the state transition to initial state and input perturbations
\cite{Sontag1998,Khalil2002Nonlinear}.

Assume
\[
\|A_t\|_2\le\rho<1,
\qquad
\|B_t\|_2\le M_B
\]
for all \(t\). At differentiability points,
\[
D_{H_t}H_{t+1}
=
D_tA_t,
\qquad
D_{R_t}H_{t+1}
=
D_tB_t,
\]
where \(D_t\) is diagonal and \(\|D_t\|_2\le1\). The same bounds hold for admissible Clarke generalized Jacobians when \(\varphi\) is piecewise affine with slopes in \([0,1]\).

\begin{proposition}[State and input sensitivity bounds]
\label{prop:state-space-sensitivity-bound}
Under the bounds above,
\[
\|D_{H_0}H_t\|_2
\le
\rho^t .
\]
For \(0\le s<t\),
\[
\|D_{R_s}H_t\|_2
\le
M_B\rho^{t-s-1}.
\]
Consequently,
\[
\alpha_H^{\mathrm{term}}
:=
\limsup_{t\to\infty}
\frac1t\log\|D_{H_0}H_t\|_2
\le
\log\rho<0
\]
and
\[
\sup_{0\le s<t}\|D_{R_s}H_t\|_2
\le
M_B .
\]
\end{proposition}

\begin{proof}
The chain rule gives
\[
D_{H_0}H_t
=
D_{t-1}A_{t-1}\cdots D_0A_0 .
\]
Since \(\|D_r\|_2\le1\) and \(\|A_r\|_2\le\rho\),
\[
\|D_{H_0}H_t\|_2
\le
\prod_{r=0}^{t-1}\|D_rA_r\|_2
\le
\rho^t .
\]
For an input perturbation at time \(s\),
\[
D_{R_s}H_t
=
D_{t-1}A_{t-1}\cdots D_{s+1}A_{s+1}D_sB_s .
\]
Therefore
\[
\|D_{R_s}H_t\|_2
\le
\left(\prod_{r=s+1}^{t-1}\|D_rA_r\|_2\right)\|D_sB_s\|_2
\le
\rho^{t-s-1}M_B .
\]
The exponent and uniform profile bounds follow directly.
\end{proof}

This example separates two stability quantities. The terminal exponent \(\alpha_H^{\mathrm{term}}\) measures forgetting of the initial hidden state. The profile bound on \(D_{R_s}H_t\) measures bounded response to later inputs or update signals. These are different stability questions, and the LSP records them separately.

Table~\ref{tab:lsp-summary} summarizes how the same stability profile records
the perturbation direction, sufficient hypotheses, and resulting bound in each
representative learning mechanism.

\begin{table}[!htbp]
\centering
\setlength{\tabcolsep}{2pt}
\renewcommand{\arraystretch}{1.25}
\caption{Representative stability mechanisms expressed through the Learning Stability Profile.}
\label{tab:lsp-summary}
%@{}p{1.50cm} | p{1.50cm} | p{4.50cm} | p{4.50cm} | p{1.90cm}@{}
\begin{tabular}{p{2.0cm} | p{2.0cm} | p{4.50cm} | p{3.50cm} | p{3.90cm}}
\toprule
\textbf{Mechanism}
&
\textbf{LSP component}
&
\textbf{Sufficient hypothesis}
&
\textbf{Representative bound}
&
\textbf{Stability meaning}
\\
\midrule

Feedforward spectral control
&
Forward sensitivity \(\sigma^x\)
&
Layerwise spectral bound
\(\|W_k\|_2\le \rho<1\)
&
\(
\|D_{X_0}X_L\|_2\le \rho^L,
\,
\alpha_x^{\mathrm{term}}\le \log\rho<0
\)
&
Terminal input perturbations decay exponentially through depth.
\\ \hline

Residual dissipativity
&
Forward sensitivity \(\sigma^x\)
&
Dissipative residual Jacobian
\(
\frac{G_k+G_k^\top}{2}\preceq -mI,
\,
\|G_k\|_2\le M_g,
\,
0<h<\frac{2m}{M_g^2}
\)
&
\(
\|D_{X_0}X_L\|_2
\le
(1-2hm+h^2M_g^2)^{L/2}
\)
&
Residual depth propagation is stable under a step-size restriction.
\\ \hline

Stochastic gradient coupling
&
Parameter and update sensitivities
\(\sigma^\theta,\sigma^u\)
&
Mean-square contraction of coupled updates:
\(
\mathbb{E}\|T_{\eta}(\theta_1,u)-T_{\eta}(\theta_2,u)\|^2
\le
\kappa_\eta^2\|\theta_1-\theta_2\|^2
\)
with \(\kappa_\eta<1\)
&
\(
\left(
\mathbb{E}\|D_{\theta^0}\theta^t\|^2
\right)^{1/2}
\le
\kappa_\eta^t,
\,
\left(
\mathbb{E}\|D_{u^s}\theta^t\|^2
\right)^{1/2}
\le
\eta L_u^G\kappa_\eta^{t-s-1}
\)
&
Initialization effects decay; update perturbations remain uniformly controlled.
\\ \hline

Piecewise-linear active sets
&
Generalized forward sensitivity
\(\sigma^{x,C}\)
&
Clarke active-set products
\(
P_L=D_{L-1}W_{L-1}\cdots D_0W_0,
\,
0\preceq D_k\preceq I
\)
&
\(
\sup_{P_L}\|P_L\|_2
\le
\prod_{k=0}^{L-1}\|W_k\|_2
\)
&
Nonsmooth ReLU-type propagation is controlled by worst-case active linearizations.
\\ \hline

Proximal or projected updates
&
Parameter sensitivity
\(\sigma^\theta\)
&
Strong convexity and smoothness:
\(
\mu I\preceq \nabla^2\mathcal{L}\preceq LI,
\,
0<\eta<\frac{2}{L+\mu}
\)
with convex \(R\)
&
\(
\|T_\eta(\theta_1)-T_\eta(\theta_2)\|
\le
q_\eta\|\theta_1-\theta_2\|,
\,
q_\eta=\max\{|1-\eta\mu|,|1-\eta L|\}
\)
&
Proximal and projected flows give contractive parameter dynamics under standard convex assumptions.
\\ \hline

Recurrent or state-space propagation
&
State and forced-input sensitivities
&
State transition bounds
\(
\|A_t\|_2\le\rho<1,
\,
\|B_t\|_2\le M_B
\)
&
\(
\|D_{H_0}H_t\|_2\le\rho^t,
\,
\|D_{R_s}H_t\|_2\le M_B\rho^{t-s-1}
\)
&
Initial hidden states are forgotten, while later inputs have bounded influence.
\\

\bottomrule
\end{tabular}
\end{table}

\section{Unified interpretation}
\label{sec:unified-interpretation}

The preceding results give a common sensitivity calculus for learning systems. The LSP records how perturbations in representations, parameters, and update variables propagate through a specified learning trajectory. The analytic and variational Lyapunov criteria give sufficient conditions under which these sensitivities are uniformly bounded or exponentially decaying. The feedforward, residual, stochastic-gradient, piecewise-linear, proximal, and recurrent examples show that the same profile separates three questions that are often mixed together, i.e., terminal decay of initial perturbations, profile-wise boundedness of forced perturbations, and the mode of stability used in deterministic, stochastic, or nonsmooth systems.

The framework should therefore be read as a stability language with explicit hypotheses, not as a universal converse theorem. Spectral bounds, dissipativity, contractive stochastic couplings, and variational Lyapunov inequalities each give different routes to controlling the same profile. When these assumptions hold globally, the resulting estimates are global. When they hold only along a trajectory or inside an invariant region, the estimates are local to that trajectory or region.

The preceding sections give a sensitivity calculus for finite-dimensional learning dynamics. The Learning Stability Profile records how perturbations in representations, parameters, and update variables propagate through a specified learning trajectory. The analytic and variational Lyapunov criteria give sufficient conditions for uniform or exponentially decaying sensitivity bounds. They do not assert an unconditional equivalence between nonpositive exponents
and dissipative energies.

\subsection{Stability as a joint property of architecture and optimization}
\label{subsec:joint-architecture-optimization}

The state variable
\[
S^t=(X_0^t,\ldots,X_L^t,\theta^t)
\]
places representations and parameters in the same dynamical system. This makes architecture and optimization part of one perturbation flow rather than two separate stability problems. The forward components of the profile measure sensitivity through representations, while the parametric and update components measure sensitivity to initialization and training perturbations.

This joint view is useful because different mechanisms act on different blocks of the same transition system. Spectral control bounds representation propagation in feedforward networks. Dissipativity and step-size restrictions control residual depth dynamics, as in explicit discretizations of dissipative systems \cite{Strikwerda2004,HaberRuthotto2017}. Contractive coupled updates control sensitivity to parameter initialization in stochastic gradient dynamics \cite{KushnerYin2003,BottouCurtisNocedal2018}. Clarke generalized Jacobians extend the same block-sensitivity calculation to locally Lipschitz maps \cite{Clarke1990Optimization,Clarke1998Nonsmooth}.

The framework therefore separates three quantities that are often conflated: terminal decay of an initial perturbation, profile-wise boundedness of forced perturbations, and the mode of stability used to interpret the bound. A negative terminal exponent gives exponential decay of the corresponding terminal sensitivity. A nonpositive profile exponent rules out exponential growth but does not imply uniform boundedness. Uniform boundedness must be proved separately.

\subsection{Relation with classical stability notions}
\label{subsec:classical-stability-relations}

The LSP does not replace classical stability theory. It records classical stability estimates in a common notation. Lyapunov stability gives transition operator bounds when a coercive energy dissipates along the linearized flow \cite{Khalil2002Nonlinear,Angeli2002Incremental}. CFL-type restrictions give admissible step sizes for explicit residual updates when the underlying Jacobian is dissipative \cite{Strikwerda2004}. Firm nonexpansiveness and averagedness give stability estimates for proximal and projected methods \cite{BauschkeCombettes2011}. Algorithmic stability controls sensitivity of a learning rule to data perturbations \cite{Bousquet2002,Hardt2016Train}.

Within the present formulation, these estimates correspond to different ways of bounding the same transition products. Spectral bounds control products of layer Jacobians. Lyapunov inequalities control products of linearized state maps. Stochastic coupling assumptions control mean-square products of random update maps. Variational assumptions control products of admissible generalized Jacobians. The mathematical content remains the estimate on the relevant transition operator; the LSP specifies which perturbation direction and which growth notion are being measured.

This also clarifies the role of nonsmooth analysis. When a learning map is differentiable, the profile is computed from classical Jacobians. When the map is locally Lipschitz, the profile is computed over admissible Clarke linearizations. At differentiability points, the Clarke generalized Jacobian reduces to the classical derivative \cite{Clarke1990Optimization}. Thus the nonsmooth formulation extends the smooth one without changing the sensitivity
object being estimated.

\subsection{Design consequences}
\label{subsec:design-consequences}

The estimates give explicit sufficient conditions for stability. For feedforward networks, layerwise spectral bounds imply
\[
\|D_{X_0}X_L\|\le \prod_{k=0}^{L-1}\|W_k\|,
\]
so a uniform bound \(\|W_k\|\le\rho<1\) gives a negative terminal forward exponent. For residual networks, the dissipativity estimate
\[
\left\|\frac{\partial X_{k+1}}{\partial X_k}\right\|_2^2
\le
1-2hm+h^2M_g^2
\]
gives a depthwise contraction whenever \(0<h<2m/M_g^2\). For stochastic gradient methods, the optimization error recursion contains the factor
\[
q_\eta=1-2\mu\eta+\eta^2(L^2+\sigma_1^2),
\]
while sensitivity to initialization requires a separate contraction assumption on coupled update maps.

These examples show why stability cannot be read from one scalar alone. A method may have bounded optimization error but unstable sensitivity to update perturbations. Conversely, a constant-step stochastic method may have a nonzero error floor while still forgetting initialization under a contractive coupling. The LSP keeps these cases distinct by reporting separate components for representation, parameter, and update sensitivity.

The framework also provides a disciplined way to state design trade-offs. Spectral normalization, residual step-size control, proximal nonexpansiveness, and stochastic coupling assumptions all reduce perturbation amplification, but they may also constrain expressivity, convergence speed, or adaptation. The profile records the stability side of this trade-off. It does not by itself settle approximation, optimization, or statistical efficiency questions.

\subsection{Robustness, generalization, and explanations}
\label{subsec:robustness-generalization-explanations}

Bounded sensitivity is closely related to robustness. If the relevant LSP component is uniformly bounded, perturbations in inputs, initialization, or update variables cannot be amplified beyond the corresponding profile bound. This gives a mechanism-level robustness statement for the specified learning trajectory.

The relation to generalization is indirect. Algorithmic stability bounds generalization by controlling sensitivity of the learned rule to changes in the training data \cite{Bousquet2002,Hardt2016Train}. The present profile can represent such sensitivity when data perturbations are encoded as parameter or update perturbations. It does not automatically imply a generalization bound without the additional assumptions required by algorithmic stability theory.

The same caution applies to explanations. Gradient-based explanation maps inherit sensitivity from network Jacobians and from the stability of the trained parameters. Instability of explanations under small perturbations has been studied through algorithmic stability measures for explanation methods \cite{Fel2022}. The LSP isolates one structural source of such instability, i.e., amplification along representation or parameter directions. It does not address all aspects of interpretability, such as semantic fidelity or domain alignment
\cite{Abusitta2024}.

\subsection{Scope and limitations}
\label{subsec:scope-limitations}

The framework is finite-dimensional and trajectory-based. Global estimates require global assumptions, such as spectral bounds, dissipativity, Lipschitz bounds, or Lyapunov inequalities holding on the relevant state region. Local assumptions yield local or trajectory-dependent estimates.

The framework also depends on the chosen stability mode. Deterministic systems are controlled through operator norms. Stochastic systems require a specified mode, such as almost sure, in expectation, in mean square, or with high probability. Nonsmooth systems require generalized derivatives or variational transition sets. These choices change the meaning of the resulting stability bound.

Several regimes are outside the present claims. The results do not give sharp bounds for arbitrary adaptive optimizers, nonstationary data streams, heavy-tailed noise without moment control, unbounded trajectories, or infinite-dimensional limits. Such settings may still be analyzed with the profile, but they require additional assumptions and separate estimates. The contribution here is structural: a common sensitivity language and Lyapunov route for stability analysis with explicit hypotheses.

\section{Conclusion}
\label{sec:conclusion}

We developed a finite-dimensional analytic and variational framework for tracking stability in learning systems through the Learning Stability Profile. The profile records sensitivity to representation perturbations, parameter initialization, and update perturbations within one coupled representation--parameter dynamics.

The main result is a Lyapunov stability criterion. In the smooth case, a coercive incremental Lyapunov energy with uniform dissipation gives explicit bounds on the linearized transition operators. In the locally Lipschitz case, the same argument controls admissible Clarke generalized transition products. These are sufficient stability criteria. The paper does not claim an unconditional converse or an equivalence between nonpositive exponents and dissipative energies.

The examples show how standard stability mechanisms fit into the profile, i.e., spectral control in feedforward networks, step-size and dissipativity control in residual networks, mean-square sensitivity bounds for stochastic gradient methods, active-set bounds for piecewise-linear networks, proximal descent, and state-space propagation. Across these cases, the profile separates terminal decay, profile-wise boundedness, and the stability mode being used.

Future work should develop continuous-time limits, local and data-dependent versions of the profile, high-probability stochastic estimates, and practical methods for estimating stability components in large-scale models. The present paper provides the finite-dimensional sensitivity calculus on which those extensions can be built.

\bibliographystyle{unsrtnat}
\bibliography{refs_new}

\end{document}
